\chapter{Implementation du chatbot intelligent (RAG)}

\section{Introduction}
Ce chapitre presente l'implementation de la composante IA, coeur de la solution proposee.

\section{Architecture conversationnelle}
Le chatbot repose sur une approche RAG qui combine :
\begin{itemize}
  \item un modele de langage pour la generation ;
  \item une base de connaissances contextuelle par hotel ;
  \item une memoire de session pour la continuite des echanges.
\end{itemize}

\section{Pipeline de traitement d'un message}
\begin{enumerate}
  \item validation et nettoyage de l'entree utilisateur ;
  \item verification des limites d'usage (rate limiting) ;
  \item chargement du contexte hotelier et de l'historique ;
  \item construction du prompt ;
  \item appel LLM et retour de la reponse (JSON ou streaming SSE) ;
  \item enregistrement d'evenements analytiques.
\end{enumerate}

\section{Multilinguisme et personnalisation}
La plateforme prend en charge plusieurs langues afin d'ameliorer l'accessibilite internationale. Les reponses sont adaptees au contexte hotelier et au profil utilisateur.

\section{Exemple de logique de traitement}
\begin{lstlisting}[caption={Pseudo-flux de traitement d'une requete chat}]
1) Validate input
2) Apply rate-limit
3) Load hotel settings + session context
4) Build RAG prompt
5) Query model
6) Stream/return answer
7) Track analytics event
\end{lstlisting}

\begin{figure}[H]
  \centering
  \begin{tikzpicture}[
    node distance=1.0cm and 0.8cm,
    box/.style={rectangle,draw,rounded corners,align=center,minimum width=2.4cm,minimum height=0.8cm},
    >={Latex}
  ]
    \node[box] (u) {Utilisateur};
    \node[box,right=of u] (api) {API Chat};
    \node[box,right=of api] (val) {Validation\\+ Rate Limit};
    \node[box,right=of val] (ctx) {Contexte RAG\\(Hotel + Session)};
    \node[box,right=of ctx] (model) {LLM Groq};
    \node[box,below=of ctx] (track) {Tracking\\Analytics};

    \draw[->] (u) -- node[above]{message} (api);
    \draw[->] (api) -- (val);
    \draw[->] (val) -- (ctx);
    \draw[->] (ctx) -- (model);
    \draw[->] (model) -- node[above]{reponse} (api);
    \draw[->] (api) -- node[above]{SSE/JSON} (u);
    \draw[->] (api) -- (track);
  \end{tikzpicture}
  \caption{Diagramme de sequence simplifie du traitement d'un message chatbot}
\end{figure}

\section{Conclusion}
L'integration du chatbot RAG apporte une assistance intelligente et contextualisee. Le chapitre suivant aborde les modules analytics et administration.
